\documentclass[12pt,a4paper]{article}
\usepackage[left=7.6667mm,right=7.3333mm,top=6.3333mm,bottom=6.3333mm,footskip=3mm]{geometry}
\usepackage{fontspec,amsmath,amssymb,booktabs,array}
\usepackage[russian]{babel}
\setmainfont{Times New Roman}
\usepackage[hidelinks]{hyperref}
\hypersetup{pdftitle={Купленный стрэддл и грид на BTC},pdfauthor={Исследование рынка Deribit}}
\setlength{\parindent}{0pt}\setlength{\parskip}{4pt}
\setlength{\abovedisplayskip}{4pt}\setlength{\belowdisplayskip}{4pt}
\setlength{\emergencystretch}{2em}
\newcommand{\head}[1]{\par\vspace{5pt}{\large\bfseries #1}\par\vspace{1pt}}
\newcommand{\sub}[1]{\par\textbf{#1.}\enspace\ignorespaces}
\newcommand{\thinrule}{\midrule}
\begin{document}
\pagestyle{plain}
\begingroup
\fontsize{11.4}{14}\selectfont
\noindent Deribit, 1 мая 2022~г. --- 14 ноября 2026~г.\par
\begin{center}\small\begin{tabular}{lrrrr}\toprule
Стратегия & 2024 & 2025 & 2026 & CAGR\\\midrule
Стрэддл + грид & 0,00 & +3,14 & +0,18 & +1,15\\
Грид, ежегодный подбор & $-0{,}12$ & $-0{,}06$ & +0,19 & +0,0016\\
Стрэддл + дельта-хедж & 0,00 & +2,18 & +0,11 & +0,79\\
Тренд за 42 дня & +29,98 & $-6{,}75$ & +21,40 & +14,40\\
Тренд + стрэддл с гридом & +14,99 & $-2{,}48$ & +11,73 & +8,17\\\bottomrule
\end{tabular}\end{center}
\smallskip
\textbf{Гипотезы и результаты.} В таблице доходность в процентах после модельных издержек. Каждый счёт стартует с 150\,000 USDC 1 января 2024 г. Годовые столбцы отсчитываются от начала года; $\mathrm{CAGR}=(E_T/E_0)^{365/d}-1$ за $d$ суток.
\begin{enumerate}
\setlength{\itemsep}{4pt}\setlength{\parsep}{0pt}\setlength{\topsep}{3pt}
\item \textbf{Покупка стрэддла.} Гипотеза: $\mathbb E[|D_T-K|\mid\text{сигнал}]>P+C$, где $P$ и $C$ --- премия и расходы на пару. На шести исполненных входах чистый стрэддл дал $+1{,}24\%$ годовых. У связки с гридом один цикл дал 3\,525 из 4\,984 USDC прибыли; без него осталось $+0{,}34\%$ годовых.
\item \textbf{Сетка.} Если возвраты к центру оплачивают инвентарь и издержки, то на общих входах $\Pi_{\text{стрэддл+грид}}-\Pi_{\text{стрэддл}}>0$. В исходном опыте задания 3 эта разность равна $-7\,763$ USDC за 36 циклов. У самостоятельного итогового грида 35 циклов дали всего $+6{,}90$ USDC. Без лучшего цикла CAGR $-0{,}0302\%$; при участии 10\% вместо 25\% и проходе и проскальзывании по 5 б.п. CAGR $-0{,}054\%$.
\item \textbf{Частичный хедж.} Гипотеза: при $H^*=-0{,}5Q\Delta$ просадка меньше, чем у чистой пары. На тех же шести входах она снизилась с $0{,}96\%$ до $0{,}68\%$ по дневным отметкам и выходам. Цена снижения риска --- 1\,953 USDC недополученной чистой прибыли. Все три бота задания с 2024 г. дали менее 4\% годовых.
\item \textbf{Тренд.} Гипотеза продолжения движения: $\mathbb E[w_t r_{t+1}]>0$, где знак $w_t$ задаёт доходность за 42 дня. С мая 2022 г. чистый CAGR составил $+16{,}68\%$ при комиссии 3,5 б.п. и $+15{,}97\%$ при 5 б.п. Без лучшей серии при исходной комиссии осталось $+9{,}65\%$; максимальная пятиминутная просадка составила $42{,}60\%$ на полной серии.
\item \textbf{Портфель 50/50.} Гипотеза: опционная часть повышает итог относительно тех же 75\,000 USDC в тренде и 75\,000 USDC в деньгах под 0\%. С 2024 г. капитал портфеля составил 187\,946,98 USDC против 185\,366,18 USDC у такого контроля. Вклад опционной части $+2\,580{,}79$ USDC; без лучшей серии одной из частей CAGR портфеля $+2{,}46\%$.
\end{enumerate}
\sub{Следующая проверка} В исходном опыте за первый час после сигнала покупки обеих ног видны в 17 из 36 циклов. Лишь в двух по обеим ногам набрался объём в пределах лимита. Неизвестно, получили бы наши заявки исполнение в очереди. Зафиксируем календарь и сигнал; на новых входах запишем синхронные bid/ask, глубину, очередь, задержку и время с одной открытой ногой. Проверим падение IV до расчёта, скачки базиса BTC/USD--USDC, расхождение с TWAP у страйка, спред при ролле и доступность маржи при сбое биржи. Критерий продолжения: положительный чистый вклад грида против чистой пары на новых общих входах и прибыль связки при худшем исполнении. При отдельном старте в январе 2025 г. исходные автономные связка, грид и частичный хедж дали $-4{,}47\%$, $-2{,}09\%$ и $-5{,}96\%$ годовых. Для вывода о 16\% у связки нужна новая серия сделок.
\par\smallskip{\small Данные и спецификации: \href{https://docs.deribit.com/api-reference/market-data/public-get_last_trades_by_currency_and_time}{Deribit, история сделок}; \href{https://support.deribit.com/hc/en-us/articles/31424939096093-Inverse-Options}{инверсные опционы}; \href{https://support.deribit.com/hc/en-us/articles/25944746248989-Fees}{тарифы}; \href{https://docs.tardis.dev/historical-data-details/deribit}{Tardis, котировки}.}
\endgroup
\newpage

\begingroup
\fontsize{10.25}{11.55}\selectfont\setlength{\leftmargini}{1.55em}
\setlength{\parskip}{1pt}\setlength{\abovedisplayskip}{2pt}\setlength{\belowdisplayskip}{2pt}
\textbf{\large Правила итоговых ботов и статистика входов}\\[2pt]01.01.2024--14.11.2026, UTC.\par
$E$ --- капитал в USDC, $S$ --- индекс BTC, $K$ --- страйк, $D_T$ --- цена расчёта.
$[x]_a=a\lfloor x/a\rfloor$. Precision $=TP/(TP+FP)$, recall $=TP/(TP+FN)$.
Одно исполнение пакета STRD считается одной сделкой; экспирация --- денежный расчёт. Метрики сигнала STRD общие для пп. 1, 2, 4 и опционной части п. 6; метрики тренда --- для пп. 5 и 6.
\begin{enumerate}
\setlength{\itemsep}{2pt}\setlength{\parsep}{1pt}\setlength{\topsep}{2pt}
\item \textbf{Чистый стрэддл.} Каждый час при свободном счёте выбирается пакет call + put одного $K,T$: $|K/S-1|\le3\%$, до расчёта от 6 часов до 7 суток. Берётся максимум объёма за 24 часа среди пакетов с доступной IV; при равенстве --- ближайший страйк, затем ближайшая экспирация. $M=C_{BS}+P_{BS}$ по IV каждой ноги при нулевой ставке. Для $j=1,\dots,13$ недель и уже завершённых аналогов
\[
 X_j=|S_t D_{T-7j}/S_{t-7j}-K|,\;
 \bar X=\frac{\sum_j X_j-\max_j X_j}{n-1},\ n\ge8,\;
 B=\big[\min(0{,}975M,0{,}90\bar X)/S_t\big]_{0{,}0001}.
\]
Для ноги $b=\pm1$ (call/put) $d_1=\ln(S/K)/(\sigma\sqrt\tau)+\sigma\sqrt\tau/2$, $d_2=d_1-\sigma\sqrt\tau$, цена $b[S\Phi(bd_1)-K\Phi(bd_2)]$, $\Delta=b\Phi(bd_1)$, $\Gamma=\phi(d_1)/(S\sigma\sqrt\tau)$, $\tau$ в годах. С 1-й минуты до конца часа покупка по $B$ при первой встречной продаже $p\le B-0{,}0001$ BTC с доступным лотом.
$Q=[\min(E/S_t,E/S_{\rm fill},0{,}08E/c,0{,}25V)]_{0{,}1}$ BTC, где $V$ --- объём принта. Стоимость пары $c=1{,}001\sum_i(a_i+\min(0{,}0003S_{\rm fill},0{,}125a_i))$, $a_i=S_{\rm fill}(p_i+(B-p)/2)$ для ног $i=C,P$. Удержание до выплаты $Q|D_T-K|$. \textbf{6 сделок.}
На 99 встречных продажах от 0,4 BTC сигнал --- покупка, событие --- выплата выше цены принта: $|D_T-K|>pS_{\rm fill}$. Precision $5/6=83{,}33\%$, recall $5/32=15{,}625\%$.
\item \textbf{Стрэддл + грид.} Покупка и расчёт по п. 1. С первой полной минуты после покупки цена-якорь $F$ берётся по последнему закрытию, затем каждые 5 минут. $h=F\max(0{,}0015;\,0{,}35\sqrt{\sum_{i=1}^{12}r_{i,5}^2})$, где $r$ --- лог-доходности перпетуала.
\par Объём уровня $u=[\min(0{,}10E/(3F),0{,}75Q/3,0{,}5Q\Gamma h)]_{0{,}001}$, $\Gamma=\Gamma_C+\Gamma_P$. При позиции $nu$ центр $A=F-0{,}5nh$; покупка по $\lfloor A-h\rfloor$, продажа по $\lceil A+h\rceil$, $|n|\le3$. Нужны проход 2 б.п. и объём свечи $\ge4u$; двусторонний проход пропускается. За 30 минут до $T$ начинается закрытие инвентаря. \textbf{369 сделок}: 6 пакетов + 363 исполнения перпетуала.
\item \textbf{Самостоятельный грид.} С 1-й минуты часа $F=P_{\rm inv}\,\mathrm{EWMA}_{1/60}(P_{\rm lin}/P_{\rm inv})$ по завершённым минутам с торговлей обоими контрактами; $P$ --- предыдущая цена линейного перпетуала. Уровень $u=[0{,}10E/(2F_0)]_{0{,}001}$, максимум $2u$. В 2024 г. при свежей сделке и $P/F-1<-0{,}0025$ покупка по $P(1-0{,}0015)$; при $P/F-1>0{,}0025$ зеркальная продажа. Добор ещё через 15 б.п. от первого исполнения при сохранении расхождения. В 2025--2026 гг. заявки $F(1\pm0{,}012)$, добор $F(1\pm0{,}024)$. Выход к $F$ или закрытие с 20-й минуты (2024--2025), 50-й (2026). Лимиты наружу до 1 USDC, проход 2 б.п., участие до 25\% встречного объёма. \textbf{71 сделка, 35 циклов.} Сигнал --- первая покупка в цикле, событие --- последняя цена исполнения выше первой до проскальзывания. Precision $9/17=52{,}94\%$, recall $9/18=50\%$.
\item \textbf{Стрэддл + дельта-хедж.} Пакет по п. 1. Каждую минуту цель $H^*=-0{,}5Q(\Delta_C+\Delta_P)$, полоса $H^*\pm0{,}15Q$. Ниже полосы покупка до нижней границы, выше --- продажа до верхней; округление внутрь до 0,001 BTC, исполнение со следующей минуты. С $T-30$ минут цель ноль. \textbf{280 сделок}: 6 пакетов + 274 исполнения перпетуала.
\item \textbf{Тренд за 42 дня.} В 00:00 $r_t=\ln(P_t/P_{t-1})$, $v_{k,t}=2^{-1/k}v_{k,t-1}+(1-2^{-1/k})r_t^2$, $k\in\{7,30\}$; $s=\max(0{,}20,\sqrt{365(0{,}65v_7+0{,}35v_{30})})$.
$w=\min(1,0{,}45/s)\,\mathrm{clip}(\ln(P_t/P_{t-42})/(s\sqrt{42/365}),-1,1)$.
С 00:05 на первом доступном открытии $P_o$ цель $H=\mathrm{sign}(w)[|w|E/P_o]_{0{,}0001}$. Покупка при росте цели относительно позиции, продажа при снижении, закрытие в конце теста. \textbf{1036 сделок.} На 1048 дневных сигналах прогноз роста $w>0$; факт --- следующая цена исполнения дневной цели выше текущей (для последнего дня --- конечное закрытие). Precision $305/591=51{,}61\%$, recall $305/528=57{,}77\%$.
\item \textbf{Тренд + стрэддл с гридом.} $E_{2,0}=E_{5,0}=75\,000$, $E_t=E_{2,t}+E_{5,t}$. Сделки по пп. 2 и 5 от капитала каждой части, без переводов. \textbf{1466 сделок}: 6 пакетов, 433 исполнения сетки, 1027 трендовых.
\end{enumerate}
\textbf{Исходный протокол для сравнения на стр.\,3.} Каждую вторую пятницу с 06.05.2022 в 12:00 UTC: ближайший общий страйк в пределах 3\% от $S$, расчёт через 14 дней в 08:00. Нужны сделки обеих ног за 6 часов. Если размах восьми дневных отметок $S$ за неделю не меньше $P_0=\sum_i[M_i+\max(0{,}025M_i,0{,}00005S)]$, в 12:01 покупаем $Q=[\min(0{,}08E/D_0,E/S)]_{0{,}1}$ пар до расчёта ($D_0$ --- премия $P$ со входными расходами). Шаг $h=0{,}005S$ до 0,1 USDC; границы $K\pm Nh$, $N=\lfloor P/h\rfloor$. Уровень $q=[Q\min(0{,}5\Gamma h,0{,}75/(N+|i_0|))]_{0{,}0001}$, $i_0$ --- начальное смещение в шагах. На соседнем нижнем уровне покупка, на верхнем продажа. Пробой по минутному закрытию: выход на следующем открытии; иначе выход с 07:30 перед расчётом. У бенчмарка полной дельты цель $-Q\Delta$, порог $0{,}05Q$. Автономный грид входил при $|\sum_1^7 r_j|/\sum_1^7|r_j|\le1/3$ и держал остаток до расчётного окна; автономный хедж возвращался к $-0{,}5Q\Delta$ при отклонении $0{,}15Q$.
\endgroup

\newpage
\fontsize{11.5}{13.7}\selectfont\setlength{\parskip}{4pt}\setlength{\tabcolsep}{5pt}
\textbf{\large Исходный эксперимент задания 3: что даёт сетка}\\[2pt]
01.05.2022--14.11.2026. Все четыре конструкции используют одинаковые даты, объёмы позиций и начальный капитал 150\,000 USDC. Просадка учитывает минутные отметки. Sharpe использует дневные доходности, множитель $\sqrt{365}$ и ставку 0\%. Последний столбец показывает среднее двух худших результатов выборки (36 циклов).
\begin{center}\normalsize
\begin{tabular}{@{}lrrrrrr@{}}\toprule
Конструкция & P\&L, USDC & CAGR & Просадка & Sharpe & Побед & Худшие 5\%, USDC\\\midrule
Стрэддл + грид & $+57\,553$ & $+7{,}41\%$ & 39,07\% & 0,49 & 16/36 & $-17\,383$\\
Стрэддл & $+65\,316$ & $+8{,}28\%$ & 35,58\% & 0,54 & 16/36 & $-18\,947$\\
Грид & $-7\,763$ & $-1{,}16\%$ & 10,35\% & $-0{,}41$ & 14/36 & $-2\,508$\\
Стрэддл + дельта-хедж & $-106\,378$ & $-23{,}81\%$ & 74,51\% & $-2{,}53$ & 9/36 & $-10\,955$\\\bottomrule
\end{tabular}\end{center}
\sub{Источники прибыли} $\Pi=Q(|D_T-K|-M)+G+F-C$, где $D_T$ --- расчётный индекс, $K$ --- страйк, $M$ --- оценка пары при входе, $G$ --- валовой результат сетки, $F$ --- фандинг, $C$ --- расходы. Опционы дали $+89\,799$ USDC: дельта $+47\,125$, гамма $+326\,651$, тета $-303\,037$, вега $+21\,268$, остаток $-2\,208$. Завершённые пары сделок сетки принесли $+7\,476$, оставшийся инвентарь стоил $13\,044$, фандинг добавил $27$. Расходы составили 26\,706 USDC: спред и тик $16\,913$, опционные комиссии $5\,375$, перпетуал $2\,166$, капитал $1\,093$, конвертация $1\,161$. Суммы округлены до целых USDC.
\sub{Исполнение} Комиссия покупки каждой ноги равна $\min(0{,}0003\cdot\text{номинал},\,0{,}125\cdot\text{премия})$. Комиссия расчёта равна $\min(0{,}00015\cdot\text{номинал},\,0{,}125\cdot\text{выплата})$. Конвертация стоит 10 б.п. Перпетуал: мейкер 1,5 б.п., тейкер 5 б.п., проскальзывание 2 б.п. Исполнение лимитной заявки требует прохода 2 б.п. и достаточного минутного объёма. Двусторонний проход бот пропускает. Модель учитывает часовой фандинг. Стоимость капитала --- 5\% годовых; база начисления включает премию, комиссию покупки и резерв 10\% текущего номинала грида.
\sub{Когда окупается премия} Полный дельта-хедж даёт приближённый локальный результат $\tfrac12 Q\Gamma S^2(\sigma_{\rm RV}^2-\sigma_{\rm IV}^2)\,dt$ до издержек. Связка сохраняет направленную составляющую $(Q\Delta+x)\,dS$. Интервалы $RV/IV<0{,}85$, $[0{,}85;1)$, $[1;1{,}25)$ и $[1{,}25;\infty)$ содержат 10, 12, 7 и 7 циклов. Средний чистый P\&L соответственно $-5\,295$, $-3\,742$, $-1\,036$ и $+23\,237$ USDC. $RV$ учитывает 15-минутные доходности после входа. Интервал $1\le RV/IV<1{,}25$ дал средний убыток. Завершённая пара сделок сетки приносит около $q(h-2f_m S)$. Частое хеджирование увеличивает оборот: комиссии и проскальзывание перпетуала стоили пороговому дельта-хеджу 82\,775 USDC.
Каждому циклу присвоен режим после завершения. Результаты связки и полного хеджа соответственно $+119\,284$ и $-16\,867$ USDC в тренде, $-184\,419$ и $-111\,326$ в боковике, $+122\,688$ и $+21\,816$ при всплесках волатильности.
\sub{Три цикла и риск пробоя} Премия и выплата относятся к одной паре, P\&L --- ко всему счёту. Суммы указаны в USDC.
\begin{center}\normalsize\begin{tabular}{@{}llrrrrr@{}}\toprule
Цикл & Вход & Пар & Премия & Выплата & Грид, чист. & Итого\\\midrule
Худший & 09.08.2024 & 4,0 & 5\,309 & 67 & $+1\,722$ & $-19\,486$\\
Типичный & 01.07.2022 & 4,7 & 2\,539 & 1\,807 & $-925$ & $-4\,476$\\
Лучший & 10.03.2023 & 7,6 & 1\,729 & 8\,316 & $-1\,931$ & $+47\,904$\\\bottomrule
\end{tabular}\end{center}
Худший цикл завершился около страйка, границы сетки выдержали. В типичном цикле у верхней границы остался шорт 0,7436 BTC номиналом 16\,006 USDC. Скачок цены на $d$ перед закрытием стоит шорту $|x|d$. При стабильном базисе опционное крыло глубоко в деньгах приносит около $Qd$. Позиция $|x|\ge Q$ исчерпывает покрытие; около страйка покрытие определяет текущая $|Q\Delta|$. Лучший цикл сохранил движение цены до расчёта. Исключение этого цикла снижает CAGR связки до $+1{,}62\%$. Квартили P\&L выборки (36 циклов): $-9\,011$, $-4\,615$, $+8\,332$ USDC.
\sub{Чувствительность} Каждый тест меняет один параметр. Сравнение сроков начинается 01.08.2025.
\begin{center}\normalsize\begin{tabular}{@{}lll@{}}\toprule
Параметр & Значения & CAGR связки, \%\\\midrule
Шаг, б.п. & 25 / 50 / 100 & 7,12 / 7,41 / 7,07\\
Объём, доля $\Gamma h$ & 0,25 / 0,50 / 1,00 & 7,66 / 7,41 / 6,46\\
Полуширина, доля премии & 0,50 / 0,75 / 1,00 & 7,88 / 7,31 / 7,41\\
Срок, дней & 7 / 14 / 21 / 28 & $-11{,}62$ / $-0{,}09$ / $-1{,}52$ / 0,32\\\bottomrule
\end{tabular}\end{center}
Календарный сдвиг на неделю дал $-6{,}46\%$ годовых. Удвоение полуспреда и стоимости конвертации вместе с проходом и проскальзыванием по 5 б.п. снизило CAGR до $+3{,}32\%$. Блочный бутстрэп (блок --- четыре двухнедельных слота) дал 95\%-й интервал $[-10{,}07;30{,}59]\%$. Расчёт интервала не учитывает подбор правил.
\end{document}
